\begin{frame}{Discussion}
  \section{Discussion}
 
Both models show western boundary currents due to mass conservation, but only Munk theory shows a western boundary current. Why?:



\begin{columns}[T] 
    \begin{column}{0.6\textwidth}

    \begin{itemize}
        \item The Sverdrup theory is linear, only one B.C. $\Psi\left(x_\text{East}\right)=0$.
        \item The Munk theory is non-linear, four B.C.s: $\Psi\left(x_\text{East}\right)=\Psi\left(x_\text{West}\right)=v\left(x_\text{East}\right)=v\left(x_\text{West}\right)=0$.
        \item Frictional force is opposite sign to the motion itself, needed to balance Sverdrup flow. 
        \item The Munk viscosity creates a boundary layer where the frictional force can take care of mass conservation.
    \end{itemize}

    \end{column}
    \begin{column}{0.4\textwidth} 
        \vspace{1cm}
      \begin{empheq}{equation*}
            A\nabla^4 \Psi -\beta \frac{\partial \Psi}{\partial x} = \frac{1}{\rho}\nabla\times\vec{\tau}
       \end{empheq} 
    \end{column}
  \end{columns}

\end{frame}
